\documentclass[../../main.tex]{subfiles}
\begin{document}

\begin{center}
  {\fontsize{48pt}{58pt}\selectfont \textbf{1989}\normalsize\textbf{年度}}
\end{center}

\begin{table}[htb]
  \centering
  \caption{出題テーマと難易度}
  \begin{tabular}{|c|p{9cm}|c|}
    \hline
    \textbf{問題番号} & \textbf{テーマ} & \textbf{難易度} \\
    \hline
    第1問 & 二曲線の交点の存在条件 & \\
    \hline
    第2問 & 放物線と円から定まる面積比の極限 & \\
    \hline
    第3問 & 上半平面の一次分数変換と反復 & \\
    \hline
    第4問 & 大きな整数の桁数と一の位 & \\
    \hline
    第5問 & 回転体の体積と積分公式 & \\
    \hline
    第6問 & 円環状の玉の配置と確率 & \\
    \hline
  \end{tabular}
\end{table}

\begin{table}[htb]
  \centering
  \caption{解答}
  \renewcommand{\arraystretch}{1.4}
  \begin{tabular}{|c|c|p{9cm}|}
    \hline
    \textbf{問題番号} & \textbf{小問} & \textbf{答え} \\
    \hline
    第1問 & -- & $k>2$ \\
    \hline
    第2問 & -- & 極限値は$2$ \\
    \hline
    第3問 & (1) & $f(z)$も上半平面に属する（証明問題） \\
    \cline{2-3}
    & (2) & $q=6\pm2\sqrt{5}$ \\
    \hline
    第4問 & -- & 整数部分は$200$桁，一の位は$9$ \\
    \hline
    第5問 & -- & $V=2\pi\displaystyle\int_0^1xf(x)\,dx=\pi$ \\
    \hline
    第6問 & -- & $\dfrac{(3k-n+1)(3k-n+2)}{(n+2)(n+1)}$ \\
    \hline
  \end{tabular}
\end{table}

\paragraph{試験形態}
6問．試験時間は要確認．

\end{document}
